\documentclass[leqno]{jsarticle}
\usepackage{amssymb}
\usepackage{amsthm}
\usepackage{amsmath}
\usepackage{comment}
	%可換図式用
		%\usepackage[all]{xy}
	%重くなるので使わいときはコメント化しておく。
%定理環境 thmはあまり使わずpropをなるべく使う。
%主張は基本的に証明の中で使い、そのため番号を付けない。
%注意環境を付けてみた。
\newtheorem{thm}{定理}
\newtheorem{prop}[thm]{命題}
\newtheorem{cor}[thm]{系}
\newtheorem{lem}[thm]{補題}
\newtheorem{df}[thm]{定義}
\newtheorem{exam}[thm]{例}
\newtheorem{fact}[thm]{事実}
\newtheorem*{claim}{主張}
\newtheorem*{cau}{注意}
\renewcommand{\proofname}{{\bf 証明}}
%矢印たち。
%\toは\rightarrowと同じ。\getsは\leftarrowと同じ。同値は\iff
%上向き矢はuparrow逆はdownarrow
\newcommand{\To}{\Longrightarrow}
\newcommand{\Gets}{\Longleftarrow}
%黒板太字
\newcommand{\nn}{\mathbb{N}}
\newcommand{\zz}{\mathbb{Z}}
\newcommand{\qq}{\mathbb{Q}}
\newcommand{\rr}{\mathbb{R}}
\newcommand{\cc}{\mathbb{C}}
%無限大
\newcommand{\mgn}{\infty}
%ギリシャン
\newcommand{\dpst}{\displaystyle}
\newcommand{\ep}{\varepsilon}
\newcommand{\al}{\alpha}
\newcommand{\be}{\beta}
\newcommand{\de}{\delta}
\newcommand{\la}{\lambda}
\newcommand{\La}{\Lambda}
\newcommand{\ga}{\gamma}
\newcommand{\lil}{\la\in \La}
%商集合
\newcommand{\qt}[2]{#1/\! #2}
%enumerate環境
\renewcommand{\labelenumi}{{\rm (\arabic{enumi})}}
%以降alignじゃなくてenumerate を使え。
%なんか演算
\newcommand{\card}{\text{{\rm card}}}
\newcommand{\supp}{\text{{\rm supp}}}
%なんか演算２
\newcommand{\bd}{\text{{\rm Bdry}}}
\newcommand{\cl}{\text{{\rm CL}}}
\newcommand{\nai}{\text{{\rm Int}}}
%差集合は\setminusで統一する。等号の否定は\neq
%真部分集合は\subsetneqq　偏微分は\partial
%ディスプレイスタイル。\dfracでディスプレイ分数
%空集合は\Oを使う！！！！！！！！！！！！

\title{電波通信}
\author{電波通信}
\date{\today}
			\begin{document}
\maketitle
この文書では有理数と無理数が$\rr$の中で稠密であることを証明する。以下では実数$x$に対して$\lfloor x \rfloor$で床関数を表す。すなわち$\lfloor x \rfloor$は$x$を越えない最大の整数のことである。この床関数に対して定義から
\[
\lfloor x\rfloor \le x<\lfloor x\rfloor+1
\]
という不等式が成り立ち、これを変形して
\[
x-1<\lfloor x\rfloor\le x
\]
という不等式を得る。この不等式に留意しておこう。まず有理数の稠密性から証明する。
\begin{prop}\label{1}
$\qq$は$\rr$の中で稠密である。
\end{prop}
\begin{proof}
$x\in \rr$を任意に与える。このとき自然数$n\in \nn$に対して床関数の不等式から
\[
nx-1<\lfloor nx\rfloor\le nx
\]
が分かるので、これを$n$で割ると
\[
x-\frac{1}{n}<\frac{\lfloor nx\rfloor}{n}\le x<x+\frac{1}{n}
\]
が分かる、つまり
\[
\left|x-\frac{\lfloor nx\rfloor}{n}\right|<\frac{1}{n}
\]
となる。$n$を大きくすれば$x$と$\dfrac{\lfloor nx\rfloor}{n}$はいくらでも距離を近くでき、$\lfloor nx\rfloor$が整数であることから$\dfrac{\lfloor nx\rfloor}{n}\in \qq$なので$\qq$が$\rr$の中で稠密であることが分かる。
\end{proof}

さて、次に移る前に次の簡単な事実に注意しよう。証明は読者に任せる。
\begin{fact}
有理数$r\in \qq$に対し、$r+\sqrt{2}$は無理数。
\end{fact}

上の事実と有理数の稠密性から無理数の稠密性がわかる。
\begin{prop}
$\rr\setminus \qq$は$\rr$の中で稠密。
\end{prop}
\begin{proof}
$x\in \rr$を任意に与える。このときもちろん$x-\sqrt{2}$も実数なので、これに対し命題\ref{1}を適応すると、任意の$\ep>0$に対し、$r\in \qq$が存在し、
\[
|(x-\sqrt{2})-r|=|x-(r+\sqrt{2})|<\ep
\]
となり、上の事実から$r+\sqrt{2}$は無理数であるから、定理の証明は終わる。
\end{proof}
















































			\end{document}



\begin{comment}
[可換図式]
\[\xymatrix{
&   & (2) \ar@{=>}[rd]&  &  &  & (5) \ar@{=>}[rd]&\\ 
& (1)\ar@{=>}[ru] \ar@{=>}[rd]&  &(4)\ar@{=>}[r]&(7)& (1)\ar@{=>}[ru] \ar@{=>}[rd]& & (7)&\\ 
&   & (3) \ar@{=>}[ur]&  &  &  &(6) \ar@{=>}[ur]&
}\]

[\bigin \endを使わない方法]

{\prop[aa] aaa}
で
\begin{prop}[aa]
aaa
\end{prop}
と同じ\df \thm \proof でも同様

[直和]
$\amalg$

[同値関係でよく使う波線]
\sim

[引数付きの新命令]
\newcommand{\prn}[1]{\|#1\|_p}

[番号のリセット]

\setcounter{equation}{0}


[字体の変更]

{\rm hogehoge}と使う。

\rm ローマン
\it イタリック
\sl 斜体

[supp]

\newcommand{\supp}{\text{{\rm supp}}}

[中央揃えの改行]

	\begin{eqnarray*}
		hoge1&=&hogehoge
		hoge2&=&hogehofe

	\end{eqnarray*}

&&の部分で揃えられる。


[\tag]
\begin{equation}
f(x) = ax^2 + bx + c \tag{1.2.3}
\end{equation}



[場合分け]

	\begin{cases}
		hoge1 & \text{状況１}\\
		hoge2 & \text{状況２}\\
		・
		・
		・
	\end{cases}



\end{comment}





